% ID: MAT-ALG-INT-015
% Titolo: Completare uguaglianze con potenze
% Disciplina: Matematica
% Area: algebra
% Argomento: numeri_interi
% Classe: Prima secondaria di secondo grado
% Difficolta: 3
% Tipo: completamento
% Tag: matematica, algebra, numeri_interi
% Fonte: appunti_matematica_anonimizzati
% Provenienza: materiale fornito dal docente; testo anonimizzato; soluzione redatta per Exergo

\begin{esercizio}
Completa senza cambiare i numeri già presenti:
\begin{enumerate}
\item $(-3)^5\cdot\boxed{\phantom{(-5)^5}}=(+15)^5$;
\item $(-2)^6\cdot(\boxed{\phantom{+}}4)^3=(+2)^{\boxed{\phantom{12}}}$.
\end{enumerate}
\end{esercizio}

\begin{soluzione}
Nel primo spazio inserisci $(-5)^5$: il prodotto delle basi è 15.
Nel secondo inserisci il segno $+$ e l'esponente 12:
$(-2)^6=2^6$ e $4^3=2^6$, dunque il prodotto è $2^{12}$.
\end{soluzione}
